% Pursuing Stacks (cotes 134-2 to 134-8): archive cote 134-4, site page 73 (PDF page 74), typescript page 161.
% typescript. Draft compared: ps3.tex:3747-3799.
% AI-transcribed by Claude Opus 5.5 (claude-opus-5-5, adaptive thinking, effort high), Message Batches API, batch ps2 (msgbatch_01NVwr1Yuo8qJf5rykeGBXPK), 2026-11-12; not yet checked by eye.
\scanpage{134-4/73}
\mspage{161}

\inserted{suggested,} by the case when M is a \uline{topos} and W the usual notion of weak equivalence,
but then to call X in M  "W-aspheric" or simply "aspheric" \inserted{does} correspond\deleted{s} to
the usual\inserted{(absolute)} notion of asphericity for the induced topos $M_{/X}$ , only in the
case when the \xout{fin} topos M itself is aspheric\unsure{,|.}  This is \xout{thx} so in the cases
I was most interested in (e.g. M a mode\over{h}{l}izing topos),  but if we want to
use it systematically in the general setting,  the \xout{xxxxxxxx followed xf}
term I used of "W-aspheric object" is definitely misleading.  Thus we
better change it now than never,  and use the word "W-object" instead,
and the notation M\hand{(W)}\inserted{\deleted{reserving the \deleted{\unsure{mathe}} "aspheric"}}\hand{.} As for the notion of a W-aspheric map,  in the
present case of a topos with the notion of weak equivalence,  it does cor-%
respond to the usual notion of asphericity for the induced morphism of to-%
poi
\[ M_{/X} \longrightarrow M_{/Y} \quad , \]
which is quite satisfactory.

THere is still need for caution with the \xout{wxxd} notion of W-aspheric
maps (W-aspheric objects have disappeared in the meanwhile!),  when wor-%
king in M = (Cat)  (and same thing if M = (Spaces)).  Namely when we got
a map \inserted{f:} $A \longrightarrow B$ in (Cat),  this is viewed for topological intuition as
corresponding to a morphism of topoi
\[ \overset{\hand{\wedge}}{A}{}^{\wedge} \longrightarrow B^{\wedge} \quad . \]
Now the requirement that  $f\colon A \to B$ should be $W_{(Cat)}$ -aspheric  is a lot
stronger than the asphericity of the corresponding morphism of topoi.
Indeed,  the latter just means that for any base change \xout{xf} in (Cat) of
the very particular "localization" o\over{f}{r} "induction" type,  namely
\[ B_{/b} \longrightarrow B \quad , \]
the corresponding map deduced by base change
\[ f_{/b} \;:\; A_{/b} \longrightarrow B_{/b} \]
is a weak equivalence\xout{x} (or equivalently,  that $A_{/b}$ is aspheric),  whereas
W-asphericity of f means that the same should hold for \uline{any} base change
$B' \to B$ in (Cat),  or equivalently,  that for \uline{any} such base-change,  with
B' having a final element moreover,  the corresponding \xout{map} category
$A'=A\mathrm{x}_B B'$  is aspheric.  To keep this distinction in mihd,  and because
the weaker notion is quite important and deserves a name definitely,
I will refer to this notion by saying f is \uline{weakly} \uline{$\underline{W}$}-\uline{aspheric}  (ret\xout{mu}rning
to the case of a general $\underline{W} \subset F\!\unsure{l}\,((Cat)))$  if \xout{the asphericity condition}
for any base change of the particular type $B_{/b} \to B$ above,  the correspon-%
\deleted{ding map A/b}

\margin{73}
